\section{Introduction}
\label{sec:introduction}

\PARstart{L}{enders} increasingly rely on machine learning models to decide who receives credit. Gradient-boosted
trees in particular have become a standard choice for credit scoring on tabular data, because they handle missing
values, mixed feature types and nonlinear interactions with little tuning~\cite{lessmann2015benchmark,ke2017lightgbm}.
A decision produced by such a model, however, affects a real person, and that person, the loan officer who
reviews the case and the regulator who audits the process all need more than a probability. They need to know why
the decision was made, what the applicant could do to obtain a better one, and whether the model treats groups of
applicants differently.

Post-hoc explanation methods such as SHAP~\cite{lundberg2017shap,lundberg2020tree} and
LIME~\cite{ribeiro2016lime} are now routinely attached to credit models to answer the first question. Recent work
on the Home Credit Default Risk dataset~\cite{homecredit2018} compared random forest, XGBoost and LightGBM and
reported SHAP and LIME explanations for the best model~\cite{shreya2025base}. Such studies stop at the explanation:
they show which features pushed a score up or down, but they do not tell an applicant what to change, they produce
static reports rather than a working decision process, and they do not examine whether the decisions differ across
groups of applicants. Three gaps in particular remain.

\emph{First, explanations are often not local.} LIME is meant to explain one prediction by fitting a simple model
in the neighbourhood of the applicant. When it is configured, as it often is, to perturb features across the whole
training distribution, the surrogate describes the model's global behaviour and gives nearly the same explanation to
everyone. In our setting, such an explanation had a rank correlation of only 0.02 with the applicant's SHAP values.

\emph{Second, recourse advice is fragile.} Counterfactual explanations~\cite{wachter2018counterfactual} and
algorithmic recourse~\cite{ustun2019recourse,karimi2022survey} go further and suggest a change that would alter the
decision, such as a smaller loan or a longer repayment period. Lenders retrain their models regularly, and an
option computed against today's model may no longer work after the next update~\cite{rawal2020shift,pawelczyk2020multiplicity}.
We found that only about a quarter of standard recourse options for rejected applicants still reached the promised
outcome on five independently retrained models. An applicant who follows such advice, for example by waiting six
months and re-applying, may be rejected again.

\emph{Third, fairness is usually treated as a model property alone.} Removing a sensitive attribute such as age from
the model is a common response to an age disparity, but other features can carry the same
information~\cite{dwork2012awareness}, and lending practice treats age partly as a question of eligibility (a minor
cannot take on a loan alone) rather than of risk.

This paper presents \emph{CreditXplain}, an explainable, counterfactual-driven decision support system for credit
risk prediction that addresses these gaps as one working system: an applicant portal, an officer dashboard and
auditable PDF reports built around a LightGBM risk model. Its contributions are as follows.

\begin{enumerate}
\item \textbf{A three-band decision process with eligibility rules.} Applications are approved automatically below a
low risk cutoff, rejected above a high one, and referred to a loan officer in between, with a written justification
required to override the model. Eligibility rules run before the model: an applicant under 18 applies with a
guardian whose income and employment are assessed, and an applicant without an income is never approved
automatically.
\item \textbf{Local explanations that agree with SHAP.} We restrict LIME to perturbations around the applicant and to
the model's most important features. Its rank correlation with SHAP rises from 0.02 to 0.64, and both explanations
are shown in plain language that reads at about an eighth-grade level.
\item \textbf{A recourse engine for credit decisions.} It searches over the loan amount, income, repayment period and
a waiting option, keeps a safety margin from the decision cutoff, rejects implausible combinations by comparison
with real applicants, respects limits set by the applicant, and offers a path in two steps, from Reject to Review and
from Review to Approve.
\item \textbf{Recourse that survives retraining.} The search requires options to work on five bootstrap-retrained
models as well as the deployed one, and is evaluated on five further retrained models that it never saw. This more
than doubles the share of options that survive retraining, from 24.4\% to 53.3\%. A hybrid that lists these options
first, and falls back to standard ones, gives 23.7\% of rejected applicants a usable option that survives retraining,
against 17.7\% for random search, 14.3\% for the standard search and 8.7\% for a baseline engine without these safeguards.
\item \textbf{An analysis of age in the model.} Removing age and the features derived from it widens the age-band
gap in automatic approval (from 0.349 to 0.377 in our experiment) and makes the model overestimate the risk of applicants under 25,
because age is largely recoverable from the remaining features. We therefore keep age as an audited feature and apply
the age policy through the eligibility rules.
\item \textbf{A complete, reproducible system.} An applicant's full path, from submission to a decision with
explanations, recourse and a PDF report, takes 158~ms on average. The code, a single script that rebuilds every model
and result, and the tests are publicly available.
\end{enumerate}

The rest of the paper is organized as follows. Section~\ref{sec:related} reviews related work.
Section~\ref{sec:method} describes the system and its components, and Section~\ref{sec:setup} the data and
evaluation protocol. Section~\ref{sec:results} reports the results, Section~\ref{sec:discussion} discusses their
limitations, and Section~\ref{sec:conclusion} concludes.
